\section{Słownik oznaczeń i krótkie podsumowanie}\label{sec:slownik}

Słownik zbiera najważniejsze oznaczenia stosowane w skrypcie. Ich znaczenie zależy od kategorii i ustalonych morfizmów.

\subsection{Kategorie i morfizmy}
\begin{center}
\renewcommand{\arraystretch}{1.15}
\begin{tabular}{@{}p{.34\linewidth}p{.60\linewidth}@{}}
\textbf{Oznaczenie} & \textbf{Znaczenie} \\
\hline
$\mathcal C,\mathcal D$ & kategorie \\
$\Ob(\mathcal C)$ & klasa obiektów kategorii \\
$\Hom_{\mathcal C}(A,B)$ & zbiór morfizmów z $A$ do $B$ (lokalna małość) \\
$g\circ f,\ \id_A$ & złożenie: najpierw $f$, następnie $g$; identyczność \\
$\mathcal C^{\op},\ f^{\op}$ & kategoria przeciwna i odwrócona strzałka \\
$A\cong B,\ f^{-1}$ & izomorfizm obiektów i odwrotność izomorfizmu \\
$\mathcal U$ & uniwersum Grothendiecka \\
$0_{A,B}$ & morfizm zerowy w kategorii z obiektem zerowym \\
\end{tabular}
\end{center}

\subsection{Przykładowe kategorie}
\begin{center}
\renewcommand{\arraystretch}{1.15}
\begin{tabular}{@{}p{.34\linewidth}p{.60\linewidth}@{}}
\textbf{Oznaczenie} & \textbf{Znaczenie} \\
\hline
$\Set,\ \Rel$ & zbiory i funkcje; zbiory i relacje \\
$\Grp,\ \Ab,\ \Mon$ & grupy, grupy abelowe i monoidy z homomorfizmami \\
$\Ring$ & pierścienie przemienne z jedynką i homomorfizmy zachowujące jedynkę \\
$\Vect_k,\ \Vect_k^{\mathrm{fd}}$ & przestrzenie wektorowe nad $k$ i pełna podkategoria przestrzeni skończenie wymiarowych \\
$\cat{Mod}_R^{\mathrm{l}}$ & lewe unitarne moduły nad pierścieniem $R$ z jedynką \\
$\Top,\ \Top_*$ & przestrzenie topologiczne i przestrzenie punktowane; odpowiednio odwzorowania ciągłe i ciągłe zachowujące punkt wyróżniony \\
$\Pos$ & zbiory częściowo uporządkowane z funkcjami monotonicznymi \\
$\mathbf 1$ & kategoria z jednym obiektem i tylko jego identycznością \\
\end{tabular}
\end{center}

\subsection{Funktory i transformacje}
\begin{center}
\renewcommand{\arraystretch}{1.15}
\begin{tabular}{@{}p{.34\linewidth}p{.60\linewidth}@{}}
\textbf{Oznaczenie} & \textbf{Znaczenie} \\
\hline
$F:\mathcal C\to\mathcal D$ & funktor kowariantny \\
$F:\mathcal C^{\op}\to\mathcal D$ & funktor kontrawariantny względem $\mathcal C$ \\
$\alpha:F\Rightarrow G$ & transformacja naturalna ze składowymi $\alpha_A:F(A)\to G(A)$ \\
$\Nat(F,G)$ & zbiór transformacji naturalnych (przy odpowiednich założeniach o wielkości) \\
$[\mathcal C,\mathcal D]$ & kategoria funktorów i transformacji naturalnych; lokalnie mała, gdy $\mathcal C$ jest mała i $\mathcal D$ lokalnie mała \\
$\beta\ast\alpha$ & składanie poziome; składanie pionowe oznaczamy $\circ$ \\
$h^A=\Hom(A,-)$ & kowariantny funktor reprezentowany przez $A$ \\
$h_A=\Hom(-,A)$ & kontrawariantny funktor reprezentowany przez $A$ \\
$\mathcal P,\ \mathcal P_+$ & funktor potęgowy kontrawariantny (przeciwobrazy) i kowariantny (obrazy); oba obejmują wszystkie podzbiory, także pusty \\
$\exists_f,\ \forall_f$ & obraz bezpośredni i zbiór punktów, których całe włókna są zawarte w zadanym podzbiorze (sekcja~\ref{sec:adjunctions}) \\
\end{tabular}
\end{center}

\subsection{Konstrukcje uniwersalne}
\begin{center}
\renewcommand{\arraystretch}{1.15}
\begin{tabular}{@{}p{.34\linewidth}p{.60\linewidth}@{}}
\textbf{Oznaczenie} & \textbf{Znaczenie} \\
\hline
$A\times B,\ A\sqcup B$ & iloczyn i koprodukt \\
$\prod_i A_i,\ \coprod_i A_i$ & iloczyn i koprodukt rodziny \\
$A\oplus B$ & suma prosta; w kategoriach addytywnych biprodukt \\
$A\times_C B,\ A\sqcup_C B$ & pullback (kwadrat kartezjański) i pushout, dla ustalonych strzałek do lub z $C$ \\
$\lim D,\ \colim D$ & granica i kogranica diagramu \\
$a\wedge b,\ a\vee b$ & kres dolny i górny w zbiorze częściowo uporządkowanym \\
$B^A,\ \ev$ & obiekt wykładniczy i ewaluacja $B^A\times A\to B$ \\
$\widetilde f$ & transpozycja $X\to B^A$ morfizmu $f:X\times A\to B$ w kontekście obiektów wykładniczych \\
$\int F$ & kategoria elementów funktora \\
$\mathcal C/B,\ B/\mathcal C$ & kategorie obiektów nad $B$ i pod $B$ \\
$(F\downarrow G)$ & kategoria przecinkowa \\
\end{tabular}
\end{center}

\subsection{Równoważności, sprzężenia i monady}
\begin{center}
\renewcommand{\arraystretch}{1.15}
\begin{tabular}{@{}p{.34\linewidth}p{.60\linewidth}@{}}
\textbf{Oznaczenie} & \textbf{Znaczenie} \\
\hline
$\mathcal C\simeq\mathcal D$ & równoważność kategorii \\
$F\cong G$ & izomorfizm naturalny funktorów \\
$F\dashv G$ & $F$ jest lewym sprzężeniem $G$ \\
$\eta,\ \varepsilon$ & jednostka $\id_{\mathcal C}\Rightarrow GF$ i kojednostka $FG\Rightarrow\id_{\mathcal D}$ \\
$(T,\eta,\mu)$ & monada: endofunktor, jednostka $\id\Rightarrow T$, mnożenie $T^2\Rightarrow T$ \\
$\mathcal C_T,\ g\star f$ & kategoria Kleisliego i jej składanie \\
$J,\ K$ & lewe i prawe sprzężenie Kleisliego, $KJ=T$ (stwierdzenie~\ref{cat:prop:kleisli-adj}) \\
$\mathcal C^T,\ (A,a)$ & kategoria Eilenberga--Moore'a i algebra monady $a:T(A)\to A$ \\
$F_T,\ U_T$ & funktory wolnej algebry i zapominania, $F_T\dashv U_T$ oraz $U_TF_T=T$ \\
$k^{(S)},\ G*H$ & wolna przestrzeń wektorowa na $S$ i iloczyn wolny grup \\
\end{tabular}
\end{center}

\subsection{Przykłady z innych działów}
\begin{center}
\renewcommand{\arraystretch}{1.15}
\begin{tabular}{@{}p{.34\linewidth}p{.60\linewidth}@{}}
\textbf{Oznaczenie} & \textbf{Znaczenie} \\
\hline
$\Ker f,\ \im f$ & jądro i obraz; w kategorii abelowej obraz to jądro kojądra \\
$\pi_1(X,x_0),\ H_n(X)$ & grupa podstawowa przestrzeni punktowanej i homologia singularna o współczynnikach całkowitych \\
$c(X),\ \pi_0(X)$ & zbiory składowych spójnych i składowych łukowej spójności \\
$\mathcal O(X),\ C(U,\mathbb R)$ & kategoria otwartych podzbiorów i zbiór funkcji ciągłych \\
$C_k(A,B)$ & funkcje ciągłe z topologią zwarto-otwartą \\
\end{tabular}
\end{center}

\subsection{Najważniejsze zależności}\label{cat:sec:summary-dependencies}

Kategoria opisuje obiekty i morfizmy (definicja~\ref{def:kategoria}), funktor zachowuje identyczności i składanie (definicja~\ref{def:funktor}), a transformacja naturalna porównuje funktory przez przemienne kwadraty (definicja~\ref{def:naturalna}). Własność uniwersalna daje istnienie i jednoznaczność odpowiedniego morfizmu.

\begin{itemize}
\item Granice i kogranice obejmują iloczyny, koprodukty, equalizatory, koequalizatory, pullbacki i pushouty (definicje~\ref{def:produkt}, \ref{def:koprodukt}, \ref{def:equalizer}, \ref{def:coeq}, \ref{def:pullback}, \ref{def:pushout}). Granica pustego diagramu jest obiektem końcowym, a kogranica początkowym (definicje~\ref{def:limit}, \ref{def:colimit}).
\item Reprezentowalność wiąże własność uniwersalną z funktorem Hom (definicja~\ref{def:representable}). Dla lokalnie małej $\mathcal C$ lemat Yonedy daje naturalne bijekcje $\Nat(h^A,F)\cong F(A)$ dla $F:\mathcal C\to\Set$ oraz $\Nat(h_A,G)\cong G(A)$ dla $G:\mathcal C^{\op}\to\Set$ (twierdzenia~\ref{tw:yoneda} i \ref{cat:thm:yoneda-contra}). Kontrawariantne funktory reprezentowalne tworzą włożenie Yonedy $A\mapsto h_A$, które jest w pełni wierne (twierdzenie~\ref{tw:yoneda-embedding}); dla funktorów kowariantnych dziedziną odpowiedniego włożenia jest $\mathcal C^{\op}$ (stwierdzenie~\ref{cat:prop:yoneda-cov-embedding}).
\item Kategoria obiektów nad bazą to szczególny przypadek kategorii przecinkowej (definicje~\ref{def:slice}, \ref{cat:def:comma}).
\item Funktor między małymi kategoriami jest równoważnością wtedy i tylko wtedy, gdy jest w pełni wierny i istotnie surjektywny na obiektach; sama pełna wierność nie wystarcza (twierdzenie~\ref{tw:equivalence}, z podanymi tam założeniami o wielkości i wyborze).
\item Sprzężenie $F\dashv G$ określa naturalne bijekcje $\Hom_{\mathcal D}(F(A),B)\cong\Hom_{\mathcal C}(A,G(B))$ (definicja~\ref{def:adjunction}). Lewe sprzężenie zachowuje istniejące małe kogranice, prawe istniejące małe granice (twierdzenie~\ref{tw:adjoints-preserve}). Konstrukcja swobodna i zapominanie to przykłady, nie definicja wszystkich sprzężeń.
\item Każde sprzężenie indukuje monadę $T=GF$ z mnożeniem $\mu_A=G(\varepsilon_{F(A)})$ (twierdzenie~\ref{cat:thm:adjunction-monad}). Każda monada na lokalnie małej kategorii pochodzi zarówno ze sprzężenia Kleisliego, jak i Eilenberga--Moore'a (stwierdzenia~\ref{cat:prop:kleisli-adj}, \ref{cat:prop:em-adj}).
\end{itemize}

\subsection{Rozróżnienia, o których warto pamiętać}\label{cat:sec:summary-pitfalls}
\begin{itemize}
\item Monomorfizm i epimorfizm definiujemy przez skracanie złożeń, nie przez injektywność i surjektywność (definicja~\ref{def:mono-epi}). Morfizm będący jednocześnie monomorfizmem i epimorfizmem nie musi być izomorfizmem (przykład~\ref{cat:ex:normed-nonabelian}).
\item Obiekt początkowy i końcowy nie muszą być tym samym obiektem. Obiekt zerowy spełnia obie własności (definicja~\ref{def:poczatkowy}).
\item Iloczyn i koprodukt mogą być różne, nawet gdy w innej kategorii realizują się przez tę samą konstrukcję. W $\Rel$ suma rozłączna jest obydwoma (przykład~\ref{cat:ex:rel-biproduct}).
\item Izomorfizm kategorii oznacza funktory wzajemnie odwrotne, a równoważność wymaga jedynie izomorfizmów naturalnych ich złożeń z identycznościami.
\item Funktor reprezentowany przez ustalony obiekt może być kowariantny lub kontrawariantny. Kierunek strzałek rozstrzyga, której postaci lematu Yonedy używamy.
\item Pełna wierność lewego sprzężenia jest równoważna odwracalności jednostki, a pełna wierność prawego sprzężenia odwracalności kojednostki (stwierdzenie~\ref{cat:prop:adj-full-faithful}). Same tożsamości trójkątne nie gwarantują odwracalności obu transformacji.
\item Istnienie iloczynów nie zapewnia istnienia obiektów wykładniczych. W $\Top$ topologia zwarto-otwarta daje obiekt wykładniczy przy wystarczającej hipotezie lokalnej zwartości i własności Hausdorffa przestrzeni $A$ (twierdzenie~\ref{cat:thm:exp-top}).
\item W kategorii abelowej jądra, kojądra i biprodukty współdziałają z addytywnością składania (definicja~\ref{cat:def:abelian-category}). Samo istnienie obiektu zerowego i skończonych iloczynów oraz koproduktów nie wystarcza do abelowości.
\end{itemize}

Najważniejsza praktyczna zasada brzmi: kiedy spotykamy nowy obiekt matematyczny, warto zapytać, \emph{jakie morfizmy do niego lub z niego wychodzą i czy można je opisać uniwersalnie}. W wielu sytuacjach taka charakterystyka jest trwalsza, prostsza i bardziej użyteczna niż konstrukcja na elementach.
